\documentclass{report}
\usepackage{amsmath,amssymb}
\setlength{\parindent}{0mm}
\setlength{\parskip}{1em}
\begin{document}
\begin{center}
\rule{6in}{1pt} \
{\large Delyan Kalchev \\
{\bf Accuracy of Least Squares Finite Element Methods for Hyperbolic Conservation Laws}}

Department of Applied Mathematics \\ University of Colorado at Boulder \\ 526 UCB \\ Boulder \\ Colorado 80309-0526
\\
{\tt delyan.kalchev@colorado.edu}\\
Christopher Leibs\\
Thomas Manteuffel\\
Steffen M�nzenmaier\end{center}

Little is currently known about the nature of the numerical error in
least-squares approximations of conservation laws. In particular, we are
interested in how accurately the dispersive properties of the continuous
equations are represented. Excessive error in dispersion relations,
introduced by the numerical approximation, leads to incorrect wave
propagation speeds. This may result in phase error, possible spurious
oscillations of non-physical nature, and incorrect representations of
wave packets. Proper understanding of dispersive errors would serve as a
foundation for improving the shock capturing capabilities of
least-squares methods. However, proper analysis is more challenging in
the context of least-squares methods than in finite-difference methods.
The purpose of this study is to reasonably quantify the capabilities of
least-squares methods to correctly capture the dispersive relations in
hyperbolic conservation laws.


\end{document}
